\documentclass[11pt]{article}
\usepackage[T1]{fontenc}
\usepackage{lmodern,microtype}
\usepackage{amsmath,amssymb,booktabs,array,enumitem}
\usepackage[letterpaper,margin=1in]{geometry}
\usepackage[colorlinks=true,urlcolor=blue!50!black,linkcolor=blue!50!black]{hyperref}
\title{Toward a metaplectic Jacquet--Langlands transfer, IV:\\Draft 1 Research and Verification Log}
\author{Prepared by ChatGPT for David Joyner}
\date{November 7, 2026}
\begin{document}
\maketitle
\section*{Scope and provenance}
This is the initial log for Part IV. The starting point was
\texttt{mjl\_part3\_draft6.tex}, its revision log, the Part II planning
memo, and the Kottwitz paper supplied in the conversation. The
full final Part I manuscript was \emph{not} available for a fresh,
line-by-line independent cocycle verification in this drafting pass.
References to Part I are therefore inputs to be rechecked when the
cross-part referee report becomes available.

The title is aligned with the rest of the series:
\emph{Toward a metaplectic Jacquet--Langlands transfer, IV: Wild
Local Correspondence and Global Trace Formula Comparison}.

\section*{New results and reductions}
\begin{enumerate}[label=\arabic*.]
\item \textbf{Integral-character-lattice reduction.}
For the unitary transform $L_D$ constructed in Part III,
integrality of coefficients in the irreducible character bases
is equivalent to permutation of those bases, with a positive sign
chosen by the identity-value normalization. The proof is elementary
unitary linear algebra, not a proof of integrality.

\item \textbf{Finite-level computation for a fixed representation.}
Because the fixed-central-character invariant functions are smooth
on compact projective quotients, a single coefficient row can in
principle be tested by a finite character calculation after
passing to an appropriate sufficiently small level. The reduction
uses the finite-expansion statement of Part III.

\item \textbf{Split diagonal-torus commutator radical.}
Under the explicitly stated standard torus cocycle
\[
\sigma_c(t,u)=(\det t,\det u)_n^c\prod_{i<j}(t_i,u_j)_n,
\]
the commutator matrix is
\[
 A=-I_r+(1+2c)J_r.
\]
When $n$ is odd, $(n,r)=1$, and $r-1+2rc\equiv0\pmod n$,
its radical on $(F^\times/F^{\times n})^r$ is the diagonal.
Consequently a regular split diagonal element is good if and only if
it belongs to $F^\times T^n$. This statement includes wild places,
but \emph{not} regular elements with nonsplit centralizers.

\item \textbf{Global arithmetic.}
The almost-everywhere local-to-global result for
$E^\times/(K^\times E^{\times n})$ from Part III is recalled with
proof, but is explicitly separated from the still-open metaplectic
local good-element and adelic measure questions.
\end{enumerate}

\section*{Verification performed}
\begin{itemize}
\item The diagonal torus commutator was calculated directly as the
ratio of the two cocycles. The matrix has entries $A_{ii}=2c$ and
$A_{ij}=1+2c$ for $i\ne j$.
\item The row sum of $A$ equals $r-1+2rc$, matching the
scalar-centrality congruence quoted in Part III. On the sum-zero
summand of $(F^\times/F^{\times n})^r$, $A$ acts as $-I$.
\item The worked numerical case $(r,n,c)=(2,3,2)$ gives
$A\equiv\left(\begin{smallmatrix}1&-1\\-1&1\end{smallmatrix}\right)\pmod3$.
\item Kottwitz (1984), \S\S5, 7--9, is used for the arithmetic
background only. Li's 2014 article is used only for the existence of
a general invariant trace formula for certain adelic covers, not for
an asserted identity in this problem.
\item Zou's 2025 simple-type paper explicitly assumes
$(p,n)=1$, so it is not a published solution of the wild-place
problem in the present draft.
\end{itemize}

\section*{Explicit unproved steps}
\begin{center}\small
\begin{tabular}{@{}p{0.39\linewidth}p{0.53\linewidth}@{}}
\toprule
Problem & What is still required\\\midrule
Part I consistency & Check the exact KP torus cocycle, inverse character conventions, and the star-map normalization\\
Wild local integrality & Prove integer character coefficients, or exhibit a counterexample to the proposed $L_D$\\
Nonsplit regular tori & Compute the commutator radical on $T_E(F)$ for every degree-$r$ finite \emph{etale} algebra $E/F$\\
Global transfer & Verify rational good-locus survival, identity and singular terms, and all Jacobian/measure conversions\\
Spectral separation & Find an actual global trace identity whose local coefficients are integers and that isolates representations\\
\bottomrule
\end{tabular}
\end{center}

\section*{Questions for the cross-part referee}
\begin{enumerate}[label=(Q\arabic*)]
\item Does the precise Part I torus cocycle restrict, up to a coboundary,
to the convention displayed above? Does its scalar-center condition
indeed match $r-1+2rc\equiv0\pmod n$?
\item Are the local $b_D=|n|_F^{-(r-1)/2}$ and
$c_D=|n|_F^{-(r^2-1)/2}$ factors correct with precisely the stated
fixed-central-character, quotient-probability Haar measures?
\item Does the Part III smoothness proof really imply finite
irreducible-character expansions in the nontrivial fixed-central-
character projective setting?
\item Are the Part II matching normalizations and the Part III trace
pairing compatible when using inverse central characters for test
functions and direct central characters for irreducibles?
\item Does the global $n$th-power theorem use all required hypotheses,
particularly $\mu_n\subset K$, $[E:K]=r$ and $(n,r)=1$?
\item Are any regular nonsplit torus good-element formulas in
Flicker--Kazhdan or Kazhdan--Patterson sufficient to eliminate the
remaining split-place task?
\end{enumerate}

\section*{Completion assessment}
Part IV, Draft 1 is a substantive beginning, not a completed proof of
metaplectic Jacquet--Langlands. The split diagonal-torus result can
be finalized after checking cocycle conventions. The wild-place
integrality problem, general local good-locus theorem, global measure
comparison, and spectral separation remain independent research
obstacles. Their difficulty should not be described as merely
editorial or a fixed number of revisions.
\end{document}
